Economic games expose strategic decisions, negotiation and long-term resource use, but a recorded action does not by itself establish choice, competence or influence. We study twenty complete MonopolyBench games containing 6,989 decisions and 8,679 response attempts, using legal menus, attempt provenance, delivered information and canonical events. Complete negotiation reconstruction identifies 920 response opportunities. Among 221 acceptance-legal offers completing a recipient color group, 34 are accepted, 113 countered and 74 rejected: own completion is insufficient for agreement. Executable audits cover all sixty eliminations and sixty-six engine defaults. Three eliminated players can pay the current debt through legal sequences; one is blocked despite sufficient nominal equity. Construction, repeated-error and asset-transfer analyses distinguish opportunity, execution and realized consequence. Selected contextual transcripts show misleading framing, explicit correction with acceptance, and commitments withdrawn under changed conditions without converting these observations into a deception rate or causal persuasion estimate. We also audit 130 supplied scenarios and run twenty separate baseline-only pilot games: eleven reference actions fail legality checks, and nineteen pilot games are censored at 1,000 decisions. All original event streams and all pilot streams replay. These results answer bounded questions about strategic execution, negotiation, economic access and measurement validity. They also specify what remains unidentified: durable learning, unfamiliar-opponent generalization, optimal policy and the relationship between general and multi-agent capability. The contribution is a source-linked empirical study and an executable procedure for separating observed behavior from stronger interpretations of agent trajectories.
